\section{Native normal-disc discovery from optimized integral cocycles}
\label{sec:cocycle-seeds}

Section~\ref{sec:diagram-exterior} supplies source-authenticated compact
exteriors. We now construct normal-surface candidates in them and return
native disc certificates when a candidate succeeds. The arithmetic idea
comes from report~31, \code{unknot\_certified\_primitives\_20261008.zip}:
minimize the number of pieces in a cocycle-dual normal surface by changing
integer vertex potentials. The maintained implementation adds primitive
cocycle discovery, a faster network solver, an independent arithmetic
checker, and an optional recognition stage. The local optimization is
polynomial. The candidate family is incomplete, with an explicit unknot
miss retained below.

\subsection{Integral cohomology in a maximal-tree gauge}

Start with the canonical finite exterior and its validated face pairings.
Choose an orientation for each global edge. The native geometry validator
already computes signed equivalences of local edges; the new producer uses
those signs to recover consistent global endpoints. The old unoriented
endpoint dictionary suffices for mod-two boundary tests but does not suffice
for this integer computation.

Choose a maximal tree in the connected one-skeleton. Any integral cocycle
can be changed by an integral vertex coboundary to vanish on the tree:
integrate its values along the unique tree paths. This gauge is unique up to
an irrelevant constant vertex potential. Put one variable on each remaining
edge and impose the oriented triangular-face equations
\[
 c([a,b])+c([b,c])-c([a,c])=0.
\]
Local signs transport each term to its global edge orientation. Coincident
global edges contribute to the same coefficient; they must not be mistaken
for distinct simplicial edges.

Sparse exact rational elimination computes the gauged kernel. The API
requires nullity one. Clear the denominators of its nonzero vector and
divide by the gcd of all its entries. The result is an integral primitive
generator: the integer kernel is saturated in the ambient edge lattice,
and a rational line meets that lattice in precisely the multiples of its
coordinate-primitive vector. This rank-one argument does not assert that
separately normalized rational basis vectors give a saturated basis in
higher rank. The implementation explicitly rejects that case.

A knot exterior has $H^1(M;\mathbb Z)\cong\mathbb Z$, so the canonical
source lies in the supported case. There are $O(T)$ variables and equations
for $T$ tetrahedra. Each original row has at most three signed contributions.
Hadamard bounds on the relevant minors give $O(T)$ bits for a primitive
kernel generator. Rational elimination uses polynomial bit cost; its
worst-case fill is quadratic and its arithmetic operation count is cubic.
The subsequent network improvement therefore does not by itself improve
this entire preprocessing bound. Work caps and cancellation also cover
elimination and the preceding manifold reconstruction.

\subsection{A normal surface without expanding its pieces}

In tetrahedron $t$, set $h_{t,0}=0$ and
$h_{t,i}=c([0,i])$. Closure makes every other edge value the corresponding
height difference. On a glued face, transported heights differ by one
integer constant. The union of half-integral affine levels therefore glues
to a cooriented normal surface. Its signed intersections with each oriented
global edge are the cocycle values, so it represents the Poincar\'e dual
of the integral class.

If the local heights, ordered by vertex, are
$a_0\le a_1\le a_2\le a_3$, the three gaps give a triangle at the lowest
vertex, a quadrilateral separating the lower pair from the upper pair,
and a triangle at the highest vertex. Zero gaps make ties harmless.
There is at most one quadrilateral type, and the number of normal pieces
is $a_3-a_0$. All seven coordinates are stored as binary integers;
neither the producer nor the native component query enumerates those pieces.

Let $v_{t,i}$ be the global vertex of local corner $i$. Changing the
cocycle by an integer potential $f$ changes the heights to
$h_{t,i}+f(v_{t,i})$. The restricted objective is
\[
 D(f)=\sum_t\left(\max_i(h_{t,i}+f(v_{t,i}))
                    -\min_i(h_{t,i}+f(v_{t,i}))\right).
\]
It measures normal pieces. It is not Euler characteristic, genus or
Thurston norm.

\subsection{The network and a linear-size independent witness}

Introduce variables $L_t,U_t$ and impose
\[
 U_t-f(v_{t,i})\ge h_{t,i},\qquad
 f(v_{t,i})-L_t\ge-h_{t,i}.
\]
Minimizing $\sum_t(U_t-L_t)$ gives the span objective. Its dual is a
three-layer transshipment with lower nodes $\ell_t$, vertex nodes $v$, and
upper nodes $u_t$. Each lower node supplies one unit and each upper node
demands one. The arcs $\ell_t\to v_{t,i}$ and $v_{t,i}\to u_t$ have
rewards $-h_{t,i}$ and $h_{t,i}$ respectively. Feasibility follows by routing
each tetrahedron through one of its own corners. Integral optimal flow
chooses one outgoing corner per lower node and one incoming corner per
upper node.

At each global vertex, match the chosen incoming and outgoing corners.
This gives $T$ quadruples $(t,i,s,j)$ with sources $t$ and targets $s$ each
appearing once and $v_{t,i}=v_{s,j}$. For every real potential,
\[
 D(f)\ \ge\ Q:=\sum_{(t,i,s,j)}(h_{s,j}-h_{t,i}),
\]
because the common vertex potential cancels in each term and the source
and target permutations sum every minimum and maximum once. An integer
potential satisfying $D(f)=Q$ proves optimality over both real and integer
potentials. The checker needs no network solver.

The maintained checker also reconstructs coordinates independently. With
$w_{ab}=|a_a-a_b|$, its triangle count at vertex $v$ is
\[
 x_v=\frac12\min_{a,b\ne v,\ a\ne b}(w_{va}+w_{vb}-w_{ab}).
\]
Write $A=w_{01}-x_0-x_1$, $B=w_{02}-x_0-x_2$ and
$C=w_{03}-x_0-x_3$. The quadrilateral counts are
$(B+C-A,A+C-B,A+B-C)/2$. For integer local heights these are nonnegative
integers and agree with the level-set surface. The checker compares the
whole vector, checks exact index types and the dual permutation, and
requires primal--dual equality. It accepts binary hexadecimal transport
without changing the interpreter's decimal-conversion limit.

\subsection{Reduced-cost augmentation improves the local bound}

Report~31 uses Bellman--Ford for each unit augmentation, with a conservative
$O(T^3)$ relaxation bound. The native solver uses a DAG initialization and
reduced-cost Dijkstra searches. Negate rewards to obtain costs, add a source
and sink with unit supply/demand arcs, and give every original internal arc
capacity $T+1$. No internal forward arc can saturate under a total flow of
$T$. The network has $O(T)$ vertices and arcs, even when heights have
thousands of bits.

Initially the usable arcs follow the order source, lower nodes, vertex
nodes, upper nodes, sink. Exact shortest distances in this DAG supply
potentials $\pi$ with nonnegative reduced costs
$c(u,v)+\pi_u-\pi_v$. In each augmentation run Dijkstra, let $d$ be its
reduced distances, and put $\Delta=d(\mathrm{sink})$. Stop as soon as the
sink is settled, prioritizing it among equal-distance choices. Update every potential
by $\min(d(v),\Delta)$, assigning the increment $\Delta$ to unreachable
or unsettled vertices. Every unsettled distance is at least $\Delta$, so
this gives the same truncated increments as a full search. This rule matters for disconnected abstract incidence
inputs; updating reachable vertices alone need not preserve all inequalities.

For any residual arc of reduced cost $c'\ge0$, the truncated increments
satisfy $a_v\le a_u+c'$. If $u$ is reachable this follows from the shortest
path inequality; otherwise $a_u=\Delta\ge a_v$. Hence all updated reduced
costs remain nonnegative. Along the augmenting path both distances are at
most $\Delta$ and each shortest-path inequality is tight, so newly introduced
reverse arcs have zero reduced cost. The invariant holds inductively.

After $T$ augmentations every original internal arc remains forward-residual,
and every used arc also has its reverse. Thus the two reduced-cost
inequalities make every used arc tight. Taking $x=-\pi$ gives feasible
$L_t=x_{\ell_t}$, $U_t=x_{u_t}$ and $f_v=x_v$. Summing tight equalities
over the integral flow proves $\sum_t(U_t-L_t)=Q$. Weak duality and
$D(f)\le\sum_t(U_t-L_t)$ then force $D(f)=Q$. The producer checks the
independent witness before returning it.

For $B$-bit heights, every original cost has magnitude at most $H<2^B$.
The source potential stays zero, and the updated sink potential is the
original cost of a simple augmenting path, bounded by $O(TH)$. Since all
increments are nonnegative and bounded by the sink increment, the other
potentials also have magnitude $O(TH)$. Arithmetic therefore needs
$O(B+\log(T+1))$ bits. A lazy binary heap gives
\[
 O(T^2\log(T+1))\ \text{arithmetic operations},\qquad
 O\!\left(T^2\log(T+1)(B+\log(T+1))\right)
\]
bit time for additions and comparisons, apart from input-label indexing
and transport. The network uses linear storage. The witness uses
$O(T\log(T+1))$ index bits and $O(T(B+\log(T+1)))$ numerical bits.
Section~\ref{sec:cocycle-blocking} subsequently batches zero-cost residual
augmentations while retaining these conservative bounds and the same checker.
These bounds are for span optimization, not all cohomology preprocessing
or complete unknot recognition.

\subsection{Connectedness is useful, but does not give minimum genus}

The report also proves a useful geometric consequence. Call a cooriented
normal surface edge coherent if all its crossings with any oriented global
edge have the same sign. Local-height surfaces have this property, which
survives deletion of component unions. Their signed edge counts form
cocycles. Admissible normal coordinates are uniquely determined by unsigned
edge counts: the kernel of the six local edge-weight equations is generated
by $(1,1,1,1,-1,-1,-1)$, and two distinct admissible vectors differing by
such a multiple would force three positive quadrilateral types in one of
them.

Suppose a nonempty component union of an optimal surface were zero in
$H_2(M,\partial M;\mathbb Z)$. Its signed cocycle would be an integral
coboundary. Deleting it preserves edge coherence and admissibility, and
the uniqueness just proved reconstructs the remaining surface from another
vertex potential in the same class. Its piece count is strictly smaller,
a contradiction. Thus there is no such zero-class union.

If $H^1(M;\mathbb Z)=\mathbb Z$ and the chosen class is primitive, each
component is nonseparating and represents a signed primitive generator.
A connected cooriented nonseparating surface is primitive because a loop
can cross it once and return through its connected complement. Oppositely
signed components would form a forbidden zero-class union. All signs
therefore agree, and a primitive total class permits only one component.
The optimized seed is connected under these hypotheses.

The initial native positive certificate retained the independent
component-level disc proof. Section~\ref{sec:cocycle-connectivity} now
checks rank-one provenance, primitivity and connectedness independently
and omits that orbit search for these restricted surfaces. In particular,
the primitive cocycle for the closure of the two-braid $(1,1,-1)$ has a
45-piece genus-one surface; span minimization reduces it to 43 pieces but
still gives genus one. That closure is an unknot. Optimizing this restricted
family is not a complete unknot test.

\subsection{Adaptive source integration and the remaining search problem}

The initial optional native stage tries the primitive cocycle surface first.
If its independent component query certifies a compressing disc, it returns
immediately. Otherwise it can run span optimization and test that second
surface. The work allowance is shared by source validation, exterior
construction, cocycle elimination, optimization, both component queries and
positive replay. The later connectedness certificates replace the two
component queries while preserving these candidate vectors and the shared
allowance. It counts cooperative guards, not bit operations or time.
The surrounding recognizer's deadline remains active throughout.

In that version every positive witness contains the exact stage-input PD, the canonical
exterior, normal coordinates and the existing normal-disc-count certificate.
A separate verifier checks the PD-to-exterior construction, replays the
surface proof and requires a strictly positive disc count. Seed discovery,
optimization and external topology engines are absent from this replay.
The certificate is bound to the PD passed to the stage; portfolio evidence
also records any preceding reductions or connected-sum factorization.

The stage follows invariant filters and an enabled group search, precedes
the optional external-engine fallback, and is disabled by default. A local
cap or an unsuccessful surface returns \code{INCONCLUSIVE}; the exact
recognizer resumes its existing path. No absent disc in one or two seeds
proves knottedness. The current successes supply native geometric proofs,
while the miss above and recorded budget exhaustions motivate further
candidate families and certified simplification. They do not establish
general quasipolynomial recognition.

\subsection{Native audits, restricted successes and retained limits}

All 507 entries of report~31's release manifest match, and all fifteen
delivered cocycle tests pass in 1.178 seconds. The maintained implementation
passes fifteen focused tests in 2.152 seconds and all 1,165 regression tests
in 215.211 seconds. A first CLI test mistakenly inspected root help instead
of the \code{recognize} subcommand; its failure log is retained, and the
final test also runs an actual CLI query and replays its returned proof.

Independent finite tests enumerate 120 small potential problems and forty
five-by-five dual assignments. Additional tests cover repeated vertices,
disconnected incidence graphs, ties, negative heights, exact schemas,
false-valued cancellation callbacks and mutations of every witness field.
Scaling an abstract instance by $2^{20000}$ leaves its network size, heap
work and augmentation count unchanged and scales the optimum exactly.
The encoded certificate still replays with producer routines disabled.
Layered-solid-torus seeds survive arbitrary local vertex and tetrahedron
relabelings. A rank-two connected-sum control is rejected; a rank-one control
with homological torsion exercises nonunit rational pivots exactly.

The pinned audit compares 200 native optimum values with the delivered
Bellman--Ford solver. Every value agrees, and the native independent checker
accepts both witnesses. The final solver makes 123,422 residual relaxation
attempts across these cases, compared with 632,778 in the delivery. These
counts measure that specific network work, not all bit operations in a
source query.

Source testing covers nineteen maintained diagrams, the circle, one named
optimized-success input, sixty valid closures retained from 180 seeded
random braids, and three larger stabilized circles: 84 cases in total.
Seventeen queries return native unknot certificates, all independently
replayed, with twelve distinct stored proofs. Sixty-two finish the restricted
candidate tests without a certified disc, and five exhaust the shared
two-million-guard allowance. No source with a known knotted verdict is
reported unknotted. Negative candidate results are never reported knotted.

The named five-crossing success is the closure of
$\sigma_1^{-1}\sigma_2\sigma_1\sigma_2^{-1}\sigma_3$ on four strands.
Its initial 194-piece surface has no compressing disc; its optimized
84-piece surface has one. The same input also occurs in the random corpus,
so the two optimized-positive records represent one distinct example.
The three-crossing unknot $(1,1,-1)$ instead remains genus one after
optimization from 45 to 43 pieces. It is an explicit limitation of the
objective, not a resource failure. Larger stabilized circles and Gordian
also expose the current shared-budget limits; those failures are retained.

For the 74 source cases with at most twelve crossings, both the raw and
optimized surfaces are independently expanded by Regina: 148 comparisons
of component counts, Euler characteristic and compressing discs all agree.
Every optimized surface is connected, as predicted by the primitive
rank-one theorem. Fourteen raw and sixteen optimized records contain discs;
these are surface controls, not additional independent knot inputs. The
maintained implementation itself uses compressed coordinates and component
orbits, not Regina or expanded sheets.

The ordinary nineteen-case audit compares the previous release, the new
code with the option disabled, and the new code with it enabled, using the
established compressed-group portfolio. All verdicts agree. The earlier
stages close every case, so there are zero native-seed attempts in that
portfolio. This evidence supports compatibility, not broader automatic
recognition coverage.

The first published native version explored full Dijkstra regions. Its
complete-call timing records showed little ordinary-height improvement.
The \path{cocycle-seed-before-sink-optimize} files retain these measurements.
Stopping at the
settled sink preserves the proof and reduces source-budget exhaustions from
twelve to five with the same seventeen positive results. The final audit
took 80.071 seconds while regression tests were also running; its elapsed
time is not used as a benchmark. Both suites and audits precede the serial
measurements below.

\subsubsection{Complete seed-optimization calls}

The final measurements run serially after the tests, audit and publication
of the runtime change. Each experiment has five shuffled paired rounds,
old and new A/A controls, and one excluded warmup per arm and input. All
580 measured calls and 116 warmups complete. The source pins are identical
across the final audit and all three benchmarks.

The optimization experiment starts from a fresh stabilized-circle diagram,
constructs and independently checks its exterior, discovers its primitive
cocycle, runs the optimizer, checks the span witness, and validates the
resulting normal coordinates. The binary case first adds a global vertex
potential of amplitude $2^{4096}$, preserving the primitive class. The
baseline is the unchanged delivered solver with a native certificate adapter;
its own built-in check is retained, and the adapter also invokes the native
independent checker. The new producer already invokes that checker. The
common measurement endpoint additionally replays the returned span proof.
These are complete adapted seed-optimization calls, not isolated flow times.

\begin{center}
\small
\begin{tabular}{lrrrrr}
\toprule
Input & Old (ms) & New (ms) & Old/new & Old A/A & New A/A \\
\midrule
1 crossing & 8.582 & 7.315 & 1.187 & 1.032 & 0.987 \\
4 crossings & 72.640 & 51.040 & 1.437 & 1.000 & 1.000 \\
8 crossings & 238.285 & 152.483 & 1.557 & 0.987 & 0.998 \\
16 crossings & 895.695 & 499.241 & 1.787 & 1.005 & 0.985 \\
4 crossings, 4096 bits & 110.527 & 58.919 & 1.906 & 0.991 & 0.999 \\
\bottomrule
\end{tabular}

\end{center}

All one hundred measured calls and twenty warmups finish; driver elapsed
time is 25.414 seconds. Paired gains range from $1.187$ to $1.906$, beyond
the observed old A/A range $0.987$--$1.032$ and new A/A range
$0.985$--$1.000$. At sixteen crossings, the median falls from 895.695 ms
to 499.241 ms. The high-bit case falls from 110.527 ms to 58.919 ms.
Each implementation's certificate hash is stable across its repeats, and
the two optimum values agree. Different optimal potentials may produce
different valid witnesses.

\subsubsection{Full recognition with the native stage deliberately reached}

The second experiment disables earlier diagram, braid and invariant
shortcuts so that the source reaches the new native stage. Both arms use
the same current recognizer; only the span optimizer changes between the
delivered implementation and the native one. The raw-circle case exits
before either optimizer. The optimized-positive case returns a checked
normal-disc proof. The remaining cases continue through the exact
Khovanov fallback after the candidate queries fail. All verdicts agree with
the known inputs; no inconclusive recognition is treated as completed.

\begin{center}
\small
\begin{tabular}{lrrrrr}
\toprule
Input & Old (ms) & New (ms) & Old/new & Old A/A & New A/A \\
\midrule
Raw circle & 376.791 & 374.341 & 1.024 & 1.010 & 0.980 \\
Optimized positive & 299.693 & 288.459 & 1.034 & 0.982 & 0.959 \\
Genus-one miss & 84.344 & 85.157 & 0.986 & 1.003 & 1.002 \\
Trefoil & 92.423 & 85.683 & 1.064 & 1.057 & 0.998 \\
Figure-eight & 161.984 & 145.843 & 1.098 & 1.000 & 0.996 \\
\bottomrule
\end{tabular}

\end{center}

All one hundred measured recognitions and twenty warmups finish in
24.132 seconds of driver elapsed time. Paired ratios range from $0.986$
to $1.098$, much smaller than the optimization-stage gains. The
optimized-positive ratio $1.034$ is comparable to its new A/A fluctuation
($0.959$), and the trefoil's $1.064$ is close to its old A/A ratio $1.057$.
The raw-circle control follows the same code in both arms. These results
do not justify a broad whole-recognition speedup claim; source preparation,
component queries, proof replay and fallback can dominate the saved flow work.

\subsubsection{The ordinary compressed-group portfolio}

The last experiment compares release \code{535a914a6} with the new release
and the native option enabled, using the established compressed-group
portfolio on all nineteen ordinary cases. The earlier stages decide every
input, so there are zero native-seed attempts in all measured calls. This
is an integration and overhead control, not a measurement of native-seed
discovery or an expansion of automatic recognition coverage.

\begin{center}
\small
\begin{tabular}{lrrrrr}
\toprule
Input & Old (ms) & New (ms) & Old/new & Old A/A & New A/A \\
\midrule
survivor-00 & 15.021 & 15.090 & 0.995 & 0.994 & 0.999 \\
survivor-01 & 12.901 & 12.983 & 0.993 & 1.002 & 1.004 \\
survivor-02 & 14.435 & 14.373 & 1.035 & 1.020 & 1.000 \\
survivor-03 & 10.080 & 10.416 & 0.977 & 1.010 & 1.029 \\
mirror-03 & 22.651 & 22.667 & 1.000 & 0.928 & 0.990 \\
survivor-04 & 8.627 & 8.598 & 0.997 & 0.997 & 0.993 \\
survivor-05 & 17.494 & 17.499 & 1.002 & 1.005 & 0.997 \\
survivor-06 & 7.098 & 7.279 & 0.975 & 1.006 & 1.050 \\
survivor-07 & 18.311 & 18.548 & 0.993 & 1.003 & 1.001 \\
survivor-08 & 6.923 & 6.898 & 0.998 & 1.016 & 1.007 \\
mirror-08 & 18.854 & 19.028 & 0.989 & 1.000 & 1.005 \\
survivor-09 & 7.055 & 6.925 & 0.985 & 1.006 & 1.017 \\
survivor-10 & 6.436 & 6.354 & 1.021 & 1.020 & 0.988 \\
survivor-11 & 6.829 & 6.842 & 1.013 & 1.013 & 0.994 \\
gordian & 1223.718 & 1233.042 & 0.998 & 1.002 & 0.986 \\
Trefoil & 0.063 & 0.064 & 0.984 & 0.987 & 0.985 \\
figure\_eight & 0.068 & 0.070 & 0.972 & 0.988 & 1.016 \\
conway & 0.980 & 0.992 & 0.985 & 1.000 & 1.014 \\
kinoshita\_terasaka & 1.006 & 1.009 & 0.998 & 0.994 & 1.001 \\
\bottomrule
\end{tabular}

\end{center}

All 380 measured calls and 76 warmups finish; driver elapsed time is
34.309 seconds. Old/new paired ratios span $0.972$--$1.035$, old A/A ratios
$0.928$--$1.020$, and new A/A ratios $0.985$--$1.050$. Gordian's old and
new medians are 1.223718 and 1.233042 seconds, with paired ratio $0.998$.
No broad ordinary-workload gain or regression is established. The new
option remains disabled by default.


Reproduction uses \path{../fast/normal_orbit_research/seeds.py}, with
\code{audit}, \code{optimize}, \code{source} and \code{pipeline} modes.
All 468 runtime, driver, test, fixture and delivery-source pins are checked
before and after each run. The pipeline baseline is commit
\code{535a914a6}; the optimization baseline is the unchanged report~31
module, identified separately in the source pins. Data use the
\code{cocycle-seed-} prefix in \path{data/}. Full source proofs, incomplete
queries and the earlier measurements remain available.

